\section{Near the boundary of the simplex}
\label{sec:frontier-boundary}

So far every positive coordinate was of order $N$
(Theorem~\ref{thm:simplex-interior} needs each to be at least $\varepsilon N$).
What about bins with some small coordinates, near an edge or a vertex of the
triangle? In this section we drop that assumption. We first give an upper bound
that holds for every positive count vector (Theorem~\ref{thm:uniform-simplex-majorant}).
Then we find the scale at which a small coordinate starts to matter. It is
$N/L^2$ when at least 2 coordinates are of order $N$
(Remark~\ref{rem:crossover-scale}). Near a vertex the picture changes: the
crossing moves to odds of order $N^{-1/2}$, and the small coordinates alone
decide it (Subsection~\ref{sec:vertex-completion}).

The start is uniform, but note that $S_n(u)$ does not involve the start.
Proposition~\ref{prop:dpmf} gives
$P^\mu_N(n)\le(\max_i\mu_i)\,p^{N-1}S_n(u)$ for every start $\mu$, so every
upper bound below also holds for a general start. The basic tool compares the
Laguerre-type polynomials of the ratio form with the Bessel series $\Phi$.

\begin{lemma}[a uniform Laguerre--Bessel inequality]
\label{lem:uniform-laguerre-bessel}
For every integer $a\ge1$ and real $w>0$,
\begin{equation}\label{eq:uniform-laguerre-bessel}
0\le1-\frac{B_a(w)}{w\Phi(aw)}\le\frac w2,
\qquad
B_a(w)=\sum_{m=1}^{a}\binom{a-1}{m-1}\frac{w^m}{m!}.
\end{equation}
\end{lemma}

The proof, in Appendix~\ref{app:boundary}, compares coefficients with
Lemma~\ref{lem:clipped-product}.

\begin{theorem}[a uniform Bessel majorant on every simplex stratum]
\label{thm:uniform-simplex-majorant}
Let $n_1,\ldots,n_k$ be arbitrary positive integers with sum $N$, and let
$0<u<1$, $v=u/(1-u)$. Define
\[
\mathcal Q_n(u)=(1-u)^{N-1}v^{k-1}
\int_0^\infty e^{-t}t^k\prod_{i=1}^k\Phi(n_i vt)\,dt.
\]
Then, with no lower bound on any $n_i/N$,
\begin{equation}\label{eq:uniform-simplex-majorant}
0\le1-\frac{S_n(u)}{\mathcal Q_n(u)}
\le k(k+1)v+\frac{k}{2}v^2\left(\sum_i\sqrt{n_i}\right)^2
\le k(k+1)v+\frac{k^2}{2}Nv^2.
\end{equation}
Thus this approximation has relative error $O_k(L^2/N)$ whenever
$u=O(L/N)$, uniformly over every positive composition of $N$.
\end{theorem}

In words, $\mathcal Q_n$ is an upper bound for every bin, and it is sharp up to
a relative error $O(L^2/N)$ in the range of the crossings. Here is the idea of
the proof, which is in Appendix~\ref{app:boundary}. We apply
Lemma~\ref{lem:uniform-laguerre-bessel} to each factor of an integral form
of $S_n$. Then the relative error is at most $kv/2$ times the mean of $t$
under the probability density proportional to
$e^{-t}t^k\prod_i\Phi(n_ivt)$, and an integration by parts bounds this mean by
$2(k+1)+v(\sum_i\sqrt{n_i})^2$.

\begin{remark}[the scale $N/(\log N)^2$]\label{rem:crossover-scale}
Here is the heuristic. A small coordinate $a$ adds its own runs. Each run needs
about 2 switches, which gives a factor $u^2$, and it can go in about $N$ places.
So what matters is $aNu^2$. At the crossings of
Section~\ref{sec:simplex-thresholds} the value $Nu=T/p$ is of order $L$, and
then $aNu^2$ is of order $aL^2/N$. So when at least 2 coordinates are of order
$N$, a rare count starts to matter once it has size $\rho N/L^2$ with $\rho$
fixed.

Appendix~\ref{app:boundary} proves this. Take $k_1\ge2$ large coordinates with
fractions $\beta_i$ among themselves and $\ell$ rare counts
$a_j=\rho_jN/L^2$. With $z=Nu$,
Theorem~\ref{thm:simplex-boundary-crossover} gives $S_n(z/N)$ up to a factor
$1+O(L^{-2}+L^2/N)$. Each rare count contributes a factor
$\Phi(A^2a_jz^2/N)$, where $A=\sum_i\sqrt{\beta_i}$, and the complete
correction of relative order $1/L$ is a factor $e^{\mathcal E(z)}$ with 3
terms. The theorem also gives the crossing as the root of an implicit
equation, and as an explicit expansion in $L$, which converges slowly. Rare
counts with $a_j=o(N/L^2)$ change that expansion only by $o(1)$
(Corollary~\ref{cor:simplex-rare-invisibility}), but their number $\ell$
changes its leading and logarithmic coefficients. With only one large
coordinate the formula would have $A=1$, so the growth factor
$e^{(A^2-1)z}$ is gone. Then the crossing moves to $u$ of order $b^{-1/2}$,
where $b$ is the large count, and the scale changes. This is the case of
Subsection~\ref{sec:vertex-completion}.
\end{remark}

\subsection{Near a vertex}\label{sec:vertex-completion}

Now let one count $b$ be large and let the other positive counts
$a_1,\ldots,a_\ell$ be small compared with $b$, with $\ell\ge1$ fixed. In
this subsection and in Appendix~\ref{app:vertex} we put
$|a|=a_1+\cdots+a_\ell$. These rare counts are not the constants $a_k$ of
Section~\ref{sec:simplex-thresholds}, which do not appear here. With
$B_a$ as in Lemma~\ref{lem:uniform-laguerre-bessel}, put
\[
H(y)=\prod_{i=1}^{\ell}B_{a_i}(y),\qquad
\kappa(y)=\frac{yH'(y)}{H(y)}.
\]
The polynomial $H$ has positive coefficients, $H(0)=0$ and $H(1)\ge1$, so
$H(y_a)=1$ has exactly one positive root $y_a$, and $y_a\le1$. Write
$S_{(b,a)}$ for $S_n$ with $n=(b,a_1,\ldots,a_\ell)$. It is a polynomial in
$u$ with nonnegative coefficients, zero constant term and a positive
coefficient at $u^\ell$, so $S_{(b,a)}(u)=1$ has exactly one positive root
$u_{b,a}$.

Each rare run again needs about 2 switches and has about $b$ places, so the
rare counts enter through $bu^2$. For fixed rare counts the bounds of
Lemma~\ref{lem:vertex-randomization} in Appendix~\ref{app:vertex} give
$S_{(b,a)}(\sqrt{y/b})\to H(y)$ as $b\to\infty$. So the crossing is near
$\sqrt{y_a/b}$, and $Nu$ is of order $\sqrt N$ instead of $L$. The next
theorem shows that this leading term stays correct uniformly as long as
$|a|=o(b)$.

\begin{theorem}[all growing rare counts near a vertex]
\label{thm:vertex-uniform-root}
Fix $\ell\ge1$. Uniformly over all positive rare counts with $|a|=o(b)$, the
unique positive crossing odds satisfy
\begin{equation}\label{eq:vertex-uniform-root}
 u_{b,a}=\sqrt{\frac{y_a}{b}}
 \left(1+O_\ell\!\left(\sqrt{\frac{|a|+1}{b}}\right)\right),
 \qquad \prod_i B_{a_i}(y_a)=1.
\end{equation}
Explicitly, put
\[
 c_\ell=\ell+\sqrt\ell+1,\quad
 \varepsilon=\sqrt{\frac{|a|+1}{b}},\quad \delta=16c_\ell\varepsilon.
\]
If $\delta\le1/2$, then
\begin{equation}\label{eq:vertex-uniform-bracket}
 (1-\delta)\sqrt{y_a/b}\ \le u_{b,a}\le
 (1+\delta)\sqrt{y_a/b}.
\end{equation}
The constants are not optimized.
\end{theorem}

In words, near a vertex the crossing is decided by the polynomial $H$ of the
rare counts alone, evaluated at $bu^2$, as long as the rare counts have total
$o(b)$. The bracket~\eqref{eq:vertex-uniform-bracket} holds for all $b$ and all
positive rare counts with $\delta\le1/2$, and it
implies~\eqref{eq:vertex-uniform-root}.

The proof is in Appendix~\ref{app:vertex}. It rests on 2 facts. First,
$\log H$ is concave, and for $0<y\le1$ its elasticity $\kappa(y)$ lies
between $\max\{\ell,\frac12\sqrt{|a|y}\}$ and $\ell+\sqrt{\ell|a|y}$
(Lemma~\ref{lem:vertex-elasticity}). Second,
$S_{(b,a)}(u)=(1-u)^{|a|}\E[H(v\mathcal T)]$ exactly, where $v=u/(1-u)$ and
$\mathcal T$ is a gamma variable plus a compound binomial sum with mean
$2+(b-1)u$ (Lemma~\ref{lem:vertex-randomization}). Jensen's inequality then
gives $S_{(b,a)}(u)\ge(1-u)^{|a|}H(bu^2)$, which exceeds $1$ at
$u=(1+\delta)\sqrt{y_a/b}$. The tangent line of $\log H$ at $bu^2$ and the
moment generating function of $\mathcal T$ give an upper bound that is
below $1$ at $u=(1-\delta)\sqrt{y_a/b}$.

We have also found the next term of $u_{b,a}$ for fixed rare counts, which
for one rare count is the 2-state fixed-edge expansion of Paper~A, the growth
$u_{b,a}\sim\ell\log|a|/(2A\sqrt{|a|b})$ with $A=\sum_i\sqrt{a_i/|a|}$ when
$|a|\to\infty$ with proportions $a_i/|a|$ bounded away from $0$, and an example
with $|a|=o(b)$ where $S_{(b,a)}(u)\sim H(bu^2)$ fails although the root is
still correct (proofs omitted, checked numerically, see
Section~\ref{sec:verify}).
